\chapter{Implementation}

This chapter describes the technical implementation details of each internship workstream.

\section{Workstream 1: Credit Decisioning Platform Modernization}

\subsection{Initial State}

The initial platform had multiple utility repositories and a manual process for compiling rule updates into \texttt{.adb} artifacts. This model created delay and introduced avoidable human dependency in the release chain.

\subsection{Repository Consolidation}

As part of modernization, utility repositories were merged into a central codebase where feasible. The consolidation reduced code scattering, simplified dependency management, and improved discoverability of rule-related components.

\subsection{Automated Rule Build Pipeline}

A CI-triggered build workflow was introduced with the following logic:

\begin{enumerate}
	\item Rule changes are committed and pushed.
	\item Pipeline trigger starts automatic compilation.
	\item Rules are compiled into deployable \texttt{.adb} artifacts.
	\item Artifacts are validated and prepared for publishing.
\end{enumerate}

This transformed a manual operational step into a repeatable machine-driven process.

\subsection{Artifact Publishing and Versioning}

Compiled artifacts are published to Artifactory with controlled versioning. Versioning provides:

\begin{itemize}
	\item Historical traceability of deployed rule packages.
	\item Rollback capability in case of regression.
	\item Clear mapping between code revision and artifact revision.
\end{itemize}

\section{Workstream 2: In-house Decision Stand-in Server}

\subsection{Architecture Selection}

A prototype-ready architecture was selected:

\begin{itemize}
	\item Backend: Java Spring Boot for robust service development and integration.
	\item Frontend: React + TypeScript for maintainable user-facing interfaces and internal operations tooling.
\end{itemize}

\subsection{Operational Role}

The stand-in server is designed as a backup-first component. It can be used when third-party decision providers are unavailable, delayed, or under maintenance. This design supports business continuity while preserving a low-risk path for incremental adoption.

\subsection{Migration Strategy}

Implementation planning included a phased migration strategy:

\begin{enumerate}
	\item Enable passive readiness as contingency capability.
	\item Execute parity testing with third-party decision outcomes.
	\item Run controlled pilots in selected domains.
	\item Gradually increase dependency on internal service as confidence grows.
\end{enumerate}

\section{Workstream 3: Fraud Central Repository with Kafka}

\subsection{Problem-Driven Design}

Fraud alerts from identity theft workflows require rapid propagation across bank systems. To avoid inconsistent decisions across teams, a central repository model was designed using asynchronous event processing.

\subsection{Event Flow Implementation}

The implemented logical flow is:

\begin{enumerate}
	\item A fraud-check transaction event is received by upstream systems.
	\item Kafka producer publishes a standardized event to a topic.
	\item Spring Boot Kafka listener consumes the event.
	\item Consumer service writes/updates record in centralized repository.
	\item Post-investigation status updates are consumed and persisted.
\end{enumerate}

\subsection{Core Design Considerations}

To ensure production suitability, the implementation design includes:

\begin{itemize}
	\item Idempotent event handling to limit duplicate-processing effects.
	\item Retry and dead-letter handling for transient failures.
	\item Status lifecycle updates to maintain auditable decision history.
	\item Monitoring hooks for consumer lag and processing health.
\end{itemize}

\section{Technology Stack Summary}

\begin{table}[h]
	\centering
	\renewcommand{\arraystretch}{1.4}
	\begin{tabular}{|C{.08\textwidth}|C{.32\textwidth}|C{.46\textwidth}|}
		\hline
		\textbf{No.} & \textbf{Technology} & \textbf{Purpose} \\
		\hline
		1 & Java Spring Boot & Backend services, integration APIs, Kafka listeners \\
		\hline
		2 & React + TypeScript & Frontend interface for stand-in decision service \\
		\hline
		3 & Kafka & Asynchronous event streaming for fraud workflows \\
		\hline
		4 & Artifactory & Versioned storage of compiled \texttt{.adb} artifacts \\
		\hline
		5 & CI Pipeline & Automated compile and publish lifecycle \\
		\hline
	\end{tabular}
	\caption{Technology stack used across internship workstreams.}
\end{table}

\section{Implementation Limitations}

Despite substantial progress, the following implementation limitations remain:

\begin{itemize}
	\item End-to-end enterprise rollout requires cross-team production approvals.
	\item Stand-in service requires long-duration parity validation before critical migration.
	\item Fraud pipeline requires deeper hardening for exactly-once semantics and advanced governance.
\end{itemize}